\section{Comparison with quantum baselines, and what hardware could run}\label{sec:sota}
The question here is not whether a quantum circuit beats a classical solver; it is how our layer compares with other \emph{quantum} treatments of the same constraints, and what a device could run. On the same residual core, with the same classically reduced objective, we compare our QAOA layer with unbalanced penalisation~\cite{montanez2024unbalanced}, with the single-basis-circuit form of Fourier-based LCU~\cite{carrera2026fourier} (one trained $R_z$ angle per count constraint and a trained single-qubit mixer; not its sampling machinery) and with QAOA with an XY mixer inside every count group, the coherent circuit that the LCU approximates. Plain penalty QAOA was the weaker baseline in \cref{sec:validation} and \cref{sec:large} and is not repeated here.

\paragraph{Gate counts.} Unbalanced penalisation and the LCU form add no two-qubit gates beyond the objective: the constraint layer is single-qubit or falls on pairs the dense objective already couples. Our layer adds the preparation and the mixers, $77\,\%$ to $14\,\%$ of the objective's gates as the core grows from $48$ to $512$ variables (before routing). After routing to heavy-hex our layer has $1.9$--$2.5\times$ the CZ gates and $2.6$--$3.7\times$ the depth of a uniform-start baseline on average over six cores of $16$--$64$ qubits; on the cores below it has $1.2$--$1.9\times$ the CZ gates of the warm-started LCU and unbalanced circuits and $0.95$--$1.4\times$ those of the XY circuit. Per layer it is therefore somewhat more expensive; what it buys is that every sample is feasible.

\paragraph{Baselines, weak and strong.} In a first comparison (\cref{tab:core_cases}) Fourier-LCU and unbalanced start from the uniform state, as in common QAOA implementations. This is a weak baseline: the Fourier-LCU paper warm-starts every qubit with $R_Y(2\arcsin\sqrt{k/n})$, a product state whose probability of satisfying a cardinality constraint is $\Theta(1/\sqrt n)$. \Cref{tab:core_baselines} gives all baselines this start and trains them with $12$ restarts ($4$ at $20$ qubits) in dense simulation; the compiled circuits equal the simulation to $10^{-15}$. Everything is on cores that were chosen to fit ($16$ and $20$ free qubits; a typical interface leaves tens of free qubits, $42$--$84$ for tight constraints at $n=96$--$192$, \cref{tab:core_map}, which no method is predicted to run).

\begin{table*}[t]
\centering\footnotesize\setlength{\tabcolsep}{4pt}
\caption{\orig{compiled counts; ideal feasibility and quality from exact simulation; success probabilities from the ESP model; nothing executed} Three cases, one $p=1$ layer on the residual core, compiled to \texttt{ibm\_fez} (best of $3$--$4$ seeds). Feasible, quality ($1$ is the optimum, among feasible samples) and $P(\mathrm{opt})$ are ideal-simulation values; the Fourier-LCU and unbalanced baselines here start from the uniform state (a weak baseline; warm-started ones in \cref{tab:core_baselines}), were trained in dense simulation with few restarts, and at $n=192$ are compiled only. The cores were chosen to fit ($16$, $20$, $30$ free qubits). Grover rows: one compiled iteration with an ancilla-free and a V-chain reflection, and the iterations needed; $|F|$: feasible strings of the core.}\label{tab:core_cases}
\begin{tabular}{rrrlrrrrrr}
\toprule
$n$ & free & $|F|$ & method & CZ & depth & ESP & feas. & qual. & $P$(opt)\\
\midrule
64 & 16 & 64 & ours (QAOA) & 473 & 541 & 0.191 & 100\,\% & 0.67 & 4.18\,\%\\
 & &  & Fourier-LCU & 412 & 525 & 0.235 & 17\,\% & 0.73 & 0.96\,\%\\
 & &  & unbalanced & 412 & 525 & 0.234 & 11\,\% & 0.65 & 0.62\,\%\\
 & & & ours (Grover, 1 it., no ancilla) & 3454 & 8025 & $5.3\cdot10^{-6}$ & \multicolumn{3}{l}{16 qubits; 6 iterations}\\
 & & & ours (Grover, 1 it., V-chain) & 323 & 578 & 0.306 & \multicolumn{3}{l}{31 qubits; 6 iterations}\\
\midrule
100 & 20 & 5000 & ours (QAOA) & 669 & 607 & 0.066 & 100\,\% & 0.65 & 0.40\,\%\\
 & &  & Fourier-LCU & 413 & 307 & 0.184 & 10\,\% & 0.65 & 0.03\,\%\\
 & &  & unbalanced & 413 & 310 & 0.184 & 12\,\% & 0.52 & 0.0086\,\%\\
 & & & ours (Grover, 1 it., no ancilla) & 6270 & 14390 & $5.1\cdot10^{-11}$ & \multicolumn{3}{l}{20 qubits; 56 iterations}\\
 & & & ours (Grover, 1 it., V-chain) & 800 & 968 & 0.033 & \multicolumn{3}{l}{40 qubits; 56 iterations}\\
\midrule
192 & 30 & 303750 & ours (QAOA) & 923 & 586 & 0.037 & 100\,\% & 0.59 & 0.0021\,\%\\
 & &  & Fourier-LCU & 612 & 414 & 0.120 & -- & -- & --\\
 & &  & unbalanced & 614 & 418 & 0.119 & -- & -- & --\\
 & & & ours (Grover, 1 it., no ancilla) & 8743 & 19188 & $2.7\cdot10^{-20}$ & \multicolumn{3}{l}{30 qubits; 433 iterations}\\
 & & & ours (Grover, 1 it., V-chain) & 1279 & 1555 & $3.2\cdot10^{-3}$ & \multicolumn{3}{l}{64 qubits; 433 iterations}\\
\bottomrule
\end{tabular}
\end{table*}


\begin{table*}[t]\centering\scriptsize\setlength{\tabcolsep}{3pt}
\caption{\orig{compiled counts; ideal feasibility, quality and $P(\mathrm{opt})$ from exact dense simulation; per-shot values from the ESP and depolarising model; nothing executed} Stronger baselines on the residual core: ws = warm start: all baselines start from the warm-started product state of the Fourier-LCU paper (RY$(2\arcsin\sqrt{k/m})$ per qubit) and are trained with $12$ restarts ($p=1$). Mean $\pm$ standard deviation over instances (objective seeds $\times$ interface tuples). Feasible, quality (among feasible samples) and $P(\mathrm{opt})$: ideal. Last column: probability of the optimum per shot under the depolarising model with the ESP of the compiled circuit.}\label{tab:core_baselines}
\begin{tabular}{llrrrrrr}\toprule
core & method & CZ & ESP & feas. & quality & $P$(opt) ideal & $P$(opt) per shot (model)\\\midrule
$n=64$, $f=16$ (6 inst.) & ours ($F$) & $352\pm100$ & $0.320\pm0.117$ & $100\pm0$\,\% & $0.71\pm0.03$ & $3.68\pm0.90$\,\% & $1.272\pm0.708$\,\%\\
 & Fourier-LCU, ws & $277\pm94$ & $0.416\pm0.138$ & $18\pm4$\,\% & $0.72\pm0.11$ & $1.37\pm0.67$\,\% & $0.585\pm0.429$\,\%\\
 & unbalanced, ws & $277\pm94$ & $0.417\pm0.139$ & $70\pm0$\,\% & $0.54\pm0.05$ & $1.15\pm0.00$\,\% & $0.481\pm0.160$\,\%\\
 & XY mixer, ws & $363\pm111$ & $0.319\pm0.124$ & $15\pm0$\,\% & $0.70\pm0.03$ & $0.63\pm0.22$\,\% & $0.217\pm0.131$\,\%\\
\midrule
$n=100$, $f=20$ (4 inst.) & ours ($F$) & $555\pm0$ & $0.162\pm0.000$ & $100\pm0$\,\% & $0.75\pm0.06$ & $0.46\pm0.40$\,\% & $0.075\pm0.065$\,\%\\
 & Fourier-LCU, ws & $299\pm0$ & $0.368\pm0.000$ & $3\pm1$\,\% & $0.72\pm0.05$ & $0.12\pm0.20$\,\% & $0.044\pm0.075$\,\%\\
 & unbalanced, ws & $296\pm3$ & $0.373\pm0.004$ & $44\pm0$\,\% & $0.57\pm0.07$ & $0.00\pm0.00$\,\% & $0.002\pm0.000$\,\%\\
 & XY mixer, ws & $391\pm0$ & $0.272\pm0.001$ & $1\pm0$\,\% & $0.76\pm0.06$ & $0.01\pm0.01$\,\% & $0.003\pm0.003$\,\%\\
\bottomrule
\end{tabular}\end{table*}


\paragraph{Results.} With the uniform start the baselines give $10$--$17\,\%$ feasible samples at $16$ and $20$ qubits, and our layer's probability of the optimum per shot under the noise model is $3.5$--$16\times$ theirs (\cref{tab:core_cases}). The warm start closes much of that gap (\cref{tab:core_baselines}). In ideal simulation at $p=1$, on the $16$-qubit core (six instances), our layer gives $100\,\%$ feasible samples, quality $0.71$ among feasible samples and $P(\mathrm{opt})=3.7\,\%$; the warm-started LCU gives $18\,\%$, $0.72$ and $1.4\,\%$, unbalanced penalisation $70\,\%$ (and $96\,\%$ at $p=2$), $0.54$ and $1.2\,\%$, and the XY mixer $15\,\%$, $0.70$ and $0.6\,\%$. At $p=2$ our layer reaches $P(\mathrm{opt})=8.3\,\%$ against $6.5\,\%$ for the LCU, whose quality ($0.82$) is above ours ($0.77$) but whose results vary by an order of magnitude between instances. At $20$ qubits (four instances, with four restarts in two and eight in the other two) the baselines are weaker, with $2$--$6\,\%$ (LCU), $44\,\%$ (unbalanced) and $1\,\%$ (XY) feasible samples. Under the depolarising model, with the ESP of each compiled circuit (our layer has more CZ gates than the warm-started LCU and unbalanced circuits), the probability of the optimum per shot is $1.3\,\%$ for our layer against $0.59\,\%$, $0.48\,\%$ and $0.22\,\%$ at $16$ qubits, between $1.2$ and $3.3$ times that of the best baseline in each of the six instances. At $20$ qubits the per-shot probability of the optimum is small for every circuit ($0.01$--$0.18\,\%$ for ours) and varies by an order of magnitude between instances: our layer is the best in three of the four instances (by $7$, $22$ and $22$ times over the best baseline) and $2.5$ times worse in one, where the single-basis LCU reaches $0.17\,\%$ against our $0.07\,\%$; with this model the unbalanced circuit stays at $0.002\,\%$ because its feasible fraction is lowered by the noise. Across all ten instances our layer is best in nine, which is a small sample. The honest summary is that the feasible-by-construction layer keeps its structural advantage ($100\,\%$ feasible samples, no penalty weights to tune) and an advantage per shot that is modest at $16$ qubits and larger but erratic at $20$, and that it depends on how well the baselines are trained.

\paragraph{Noise model.} We use three levels, in increasing cost: ESP (the product of one minus the error of every compiled gate), a global depolarising model ($\mathrm{ESP}\cdot P_\mathrm{ideal}+(1-\mathrm{ESP})\cdot\mathrm{uniform}$) and, up to about $20$ qubits, a gate-level Aer simulation with the FakeFez noise model. On one $16$-qubit core the model is only slightly conservative ($18.1\,\%$ feasible shots against $21.8\pm1.7\,\%$ simulated for our layer, $600$ shots; \cref{sec:core}). Per-shot values in this section are predictions, not measurements, and apply to the models of two small cores.

\begin{table*}[t]
\centering\footnotesize\setlength{\tabcolsep}{4pt}
\caption{\orig{estimate, an extrapolation of the ESP model; no hardware} What one QAOA layer needs. Effective error per CZ $\varepsilon=-\ln(\mathrm{ESP})/\#\mathrm{CZ}$ over nine compiled circuits (median $0.0035$). ESP of our layer at $473$, $669$, $923$ CZ if the error per CZ were reduced by the factor shown; last columns: free qubits whose layer keeps ESP $\ge0.3$ and $\ge0.1$ at $\approx31$ CZ per free qubit (a dense objective gives fewer). The scaling of the error rate is an assumption.}\label{tab:core_hw}
\begin{tabular}{lrrrrr}
\toprule
error per CZ & ESP$_{16}$ & ESP$_{20}$ & ESP$_{30}$ & free qubits, ESP$\ge0.3$ & ESP$\ge0.1$\\
\midrule
today ($\varepsilon\approx0.0035$) & 0.19 & 0.09 & 0.04 & 11 & 21\\
$3\times$ better & 0.57 & 0.46 & 0.34 & 33 & 63\\
$10\times$ better & 0.85 & 0.79 & 0.72 & 110 & 211\\
\bottomrule
\end{tabular}
\end{table*}


\paragraph{Does the block grow with the number of constraints?} \orig{exact enumeration and compiled counts to $\{$CX, Rz, SX, X$\}$, all-to-all; no noise, no device} It depends on whether the constraints decompose (\path{experiments/constraint_scaling*.py}; the three parts below are the three panels of \cref{fig:constraint_scaling}). (a)~\emph{Panel (a).} For disjoint cardinality constraints the block is one Dicke state per group and the depth is that of the largest group: at $n=24$, going from $m=1$ to $12$ groups lowers it from $1761$ to $13$, and adding groups of the same size leaves it unchanged while the CZ count grows linearly. A penalty layer is cheaper per constraint ($552$ CZ and depth $135$ against $1962$ and $1761$ for one group of $24$, or $981$ CZ with the optimised Dicke block of \cref{sec:core}) but does not give feasible samples. (b)~\emph{Panel (b).} For overlapping constraints the generic automaton circuit does grow: with random constraints (five of $14$ variables, at most two) the width goes from $3$ to about $40$ and the circuit from $293$ to $1.7\cdot10^{4}$ CZ for eight constraints, while a penalty layer grows from $20$ to $116$ CZ. In that regime the construction is not competitive in gate count. (c)~\emph{Panel (c).} Interface conditioning is what restores a small block: fixing the $s$ most frequent variables classically halves the width for each variable fixed (mean CZ per branch over two instances: $1.7\cdot10^4$ at $s=0$, $1.6\cdot10^3$ at $s=3$ and $1.6\cdot10^2$ at $s=6$), at the price of $7$ and about $33$ classical branches. So the block is constant in the number of constraints only for the part that decomposes; the rest is moved to the classical interface, whose size is the real cost of the method.

\paragraph{What a device could run.} The effective error per CZ over all nine compiled circuits is $0.0035$--$0.0041$, which gives $\mathrm{ESP}\approx e^{-0.0035\,N_\mathrm{CZ}}$, and one free qubit costs about $31$ CZ in these cores (the objective couples only variables of one period). A layer with ESP $\ge0.3$ then allows about $11$ free qubits and one with ESP $\ge0.1$ about $21$ (\cref{tab:core_hw}); our layer compiles to $473$, $669$ and $923$ CZ for $16$, $20$ and $30$ free qubits, with ESP $0.19$, $0.066$ and $0.037$. Post-selection on feasibility is free for a feasible-by-construction circuit and turns noise into a loss of shots by $1/\mathrm{ESP}$ rather than a distortion: on the $16$-qubit core the quality after post-selection is $0.713$ against $0.717$ ideal. It does not catch errors that map a feasible string to another feasible string. Error cancellation would cost about $e^{4\varepsilon N_\mathrm{CZ}}$ ($7\cdot10^2$ at $473$ CZ, $4\cdot10^5$ at $923$) and is not used; dynamical decoupling, twirling and readout mitigation were not applied in the estimates.

\paragraph{The generic automaton does not run at these sizes.} With the generic automaton preparation the layer takes $2.0$, $2.9$ and $4.9$ thousand CZ gates against $464$, $581$ and $923$ for the symmetric lowering, even with measurement-based uncomputation (\cref{sec:core}). \emph{So the part specific to us, the automaton for arbitrary constraints, is not predicted to run on today's devices at these sizes; what is predicted to run is its symmetric special case.}

\paragraph{Amplitude amplification is a fault-tolerant use.} Among the QAOA variants compared here only ours admits it. One iteration costs $3.5$, $6.3$ and $8.7$ thousand CZ with an ancilla-free reflection, and $323$, $800$ and $1279$ with a Toffoli V-chain ($10.7\times$ to $6.8\times$ fewer; $31$, $40$ and $64$ qubits), with ESP per iteration $0.31$, $0.033$ and $0.003$ against the $6$, $56$ and $433$ iterations needed. Its advantage appears for large cores, where it needs $2^{4}$--$2^{6}$ times fewer iterations than Grover over all strings in the cases above and up to $2^{95}$ in the larger cores of \cref{sec:core} ($23$--$326$ free qubits); a fault-tolerant cost model ($88$--$1156$ logical qubits) is in \cref{sec:large}.

\paragraph{Combining with other methods.} A diagonal layer of any kind keeps the support, so \emph{any approximation of the objective layer preserves feasibility} and only affects quality. (This paragraph uses a second interface tuple per size, \cref{sec:core}, so its values are comparable only within it.) Without the two-qubit objective layer, our preparation and mixers compile to $66$, $260$ and $357$ CZ gates with ESP $0.80$, $0.41$ and $0.24$, against $0.20$, $0.15$ and $0.04$ with it. Replacing the objective by a single-qubit mean-field layer, as the cheapest stand-in for a Fourier-type decomposition, keeps the quality in ideal simulation ($0.741$, $0.700$, $0.622$ against $0.731$, $0.677$, $0.599$ with the full layer). The stand-in is not the LCU method, and the mean-field problem is trivial classically, so a quantum layer on top of it adds value only if it improves on the classical solution of that linearisation, which was not tested.

\paragraph{A 120-qubit experiment, specified and not run.} The estimates above concern one core at a time. To use a chip of more than $100$ qubits, we use the independence the interface gives: a portfolio with $n=120$ variables splits into twelve blocks of $10$ qubits, each with its own quadratic objective, placed on disjoint regions of $11$ physical qubits (a chip with at least $132$ usable qubits, such as the $156$-qubit \texttt{ibm\_fez}) and run as one circuit of $120$ measured qubits. Twelve circuits are specified (ours light and full, ours preparation only, a Hadamard control, the warm-started Fourier-LCU variant, unbalanced penalty and XY-mixer QAOA, trained and given the same decomposition, which is the setting most favourable to them, and, for the cost-informed preparation of \cref{sec:core}, our tilted preparation alone and with $p=1$ against the three baselines given the same classical information), together with five criteria fixed before any run (\cref{app:hw100}). \Cref{tab:hw100_pred} gives the compiled counts and the noise-model prediction: the preparation alone is predicted to give about $67\,\%$ feasible shots per block against $28\,\%$ for the warm-started unbalanced circuit, $11$--$12\,\%$ for the LCU and XY circuits and $1.7\,\%$ for the control; our $p=1$ layer gives $22\,\%$ (light) and $13\,\%$ (full), below the unbalanced circuit, because the layer costs three to four times the preparation; and the LCU and XY circuits are predicted to have a better quality among their feasible samples ($0.71$ and $0.70$ against $0.66$). So the clearest prediction is for the preparation alone, not for the $p=1$ layer. With the classical tilt (same gates, other Dicke angles) the preparation alone is predicted to give $66\,\%$ feasible shots per block and a quality of $0.79$, against $14$--$33\,\%$ and $0.83$--$0.88$ for the baselines that receive the same information: the tilt raises every circuit's quality, and ours keeps the lead in feasibility, not in quality. Blocks of ten qubits are classically trivial, so the experiment would show feasibility at a scale where global-constraint methods have no chance, not an advantage. The experiment is released as notebooks and has not been run.

\begin{table*}[t]\centering\small
\caption{\orig{compiled counts; ideal values from exact simulation; noisy values from the depolarising model; nothing executed} Predictions for the 120-qubit experiment (twelve blocks of $10$ qubits in one circuit; averages over blocks; CZ and ESP: median per block, depth: maximum). Feasible: share of shots of a block that satisfy its constraints; quality: $1-$normalised cost among feasible shots (uniform on the feasible set: about $0.5$). Not measured.}\label{tab:hw100_pred}
\begin{tabular}{lrrrrrrr}\toprule
method & CZ & depth & ESP & feas. ideal & feas. model & qual. ideal & qual. model\\\midrule
ours, $p{=}1$ light & 338 & 1008 & 0.22 & 1.000 & 0.220 & 0.67 & 0.66\\
ours, $p{=}1$ full & 447 & 1196 & 0.11 & 1.000 & 0.134 & 0.72 & 0.69\\
ours, preparation only & 96 & 449 & 0.71 & 1.000 & 0.669 & 0.50 & 0.50\\
uniform (Hadamards) & 0 & 4 & 1.00 & 0.017 & 0.017 & 0.50 & 0.50\\
Fourier-LCU, warm start & 154 & 296 & 0.47 & 0.244 & 0.122 & 0.74 & 0.71\\
unbalanced, warm start & 154 & 299 & 0.47 & 0.589 & 0.281 & 0.51 & 0.51\\
XY mixer, warm start & 194 & 342 & 0.41 & 0.239 & 0.107 & 0.73 & 0.70\\
ours, tilted preparation only & 99 & 457 & 0.69 & 1.000 & 0.660 & 0.80 & 0.79\\
ours, tilted, $p{=}1$ light & 342 & 1028 & 0.21 & 1.000 & 0.217 & 0.85 & 0.82\\
Fourier-LCU, tilted start & 154 & 298 & 0.47 & 0.375 & 0.182 & 0.91 & 0.88\\
unbalanced, tilted start & 153 & 296 & 0.47 & 0.690 & 0.330 & 0.84 & 0.83\\
XY mixer, tilted start & 192 & 344 & 0.40 & 0.326 & 0.138 & 0.91 & 0.87\\
\bottomrule\end{tabular}\end{table*}


\paragraph{Depth and cost on IBM hardware at 120 qubits.} \orig{compiled counts and an estimate; nothing executed} \Cref{tab:hw100_cost} adds the total depth and the two-qubit depth of the same circuits. The twelve blocks run in parallel, so depth does not grow with the number of blocks: our $p=1$ light layer has depth $1008$ and two-qubit depth $310$ ($4128$ CZ in all, about $344$ per block), the full layer $1196$ and $371$, the preparation alone $449$ and $149$, and the baselines $296$--$342$ and $103$--$115$ ($1838$ CZ for the LCU and unbalanced circuits, $2289$ for the XY mixer). Our layer therefore has about three times the two-qubit depth of the LCU and unbalanced circuits, and the preparation alone about $1.4\times$ with $34\%$ fewer CZ gates; the tilt changes these numbers by less than $4\%$. The scheduled duration is $13$--$44\,\mu$s per shot, so the time on the device is set by the repetition delay (default $250\,\mu$s) and the readout, not by the circuit: about $2.6$--$2.9$ s per $10^4$ shots for each circuit, $\approx11$ s for the four-circuit comparison (ours with and without the layer, LCU, unbalanced) and $\approx32$ s for all twelve. At the pay-as-you-go price listed by IBM ($\$96$ per minute of QPU time, billed per second) this is about $\$4.3$ per circuit, $\$17$ for the four-circuit comparison and $\$51$ for all twelve per $10^4$ shots, before the job overhead, which we did not estimate; a $10^5$-shot run multiplies it by ten. Whether a $156$-qubit device is available in the free plan ($10$ minutes per $28$ days) depends on the account. These are estimates from a calibration snapshot, not measurements; the total and two-qubit depths of the $16$--$30$-qubit cores of \cref{sec:core} were not recompiled.
\begin{table}[t]\centering\scriptsize\setlength{\tabcolsep}{2pt}
\caption{\orig{compiled counts on the \texttt{ibm\_fez} snapshot; time is an estimate; nothing executed} The $120$-qubit circuits (twelve blocks in one circuit): total depth, two-qubit depth, CZ gates (all blocks), scheduled duration and estimated QPU time for $10^4$ shots (duration plus the default repetition delay of $250\,\mu$s per shot; job overhead not included).}\label{tab:hw100_cost}
\begin{tabular}{lrrrrr}\toprule
circuit & depth & 2q depth & CZ & dur. ($\mu$s) & QPU s\\\midrule
ours $p{=}1$ light & 1008 & 310 & 4128 & 37 & 2.9\\
ours $p{=}1$ full & 1196 & 371 & 5466 & 44 & 2.9\\
ours, preparation & 449 & 149 & 1220 & 18 & 2.7\\
uniform & 4 & 0 & 0 & 2 & 2.5\\
Fourier-LCU, ws & 296 & 103 & 1838 & 13 & 2.6\\
unbalanced, ws & 299 & 103 & 1838 & 13 & 2.6\\
XY mixer, ws & 342 & 115 & 2289 & 15 & 2.6\\
ours tilted, prep. & 457 & 143 & 1230 & 18 & 2.7\\
ours tilted $p{=}1$ & 1028 & 319 & 4164 & 38 & 2.9\\
LCU, tilted start & 298 & 103 & 1838 & 13 & 2.6\\
unbalanced, tilted start & 296 & 103 & 1797 & 13 & 2.6\\
XY, tilted start & 344 & 115 & 2274 & 15 & 2.6\\
\midrule
all twelve circuits & & & & & 32\\
\bottomrule\end{tabular}\end{table}

